\subsection{Results on Real Radiographs}
\label{sec:realres}
Table~\ref{tab:real} summarises the real-image experiment. Without any intervention, the model already served female patients worse: discrimination was lower (AUC 0.74 vs.\ 0.84), female patients were deferred more often (+4.8\pp), and their silent-miss rate was higher (+6.9\pp; positive in 84\% of splits). This is the signal-gap pattern of Section~\ref{sec:decomp}, in which abstention does point at the disadvantaged group.

\begin{table*}[t]
\centering
\caption{Real chest radiographs (637 images, 381 patients; DenseNet-121 embeddings; 300 patient-level splits). Rates and gaps are $B$ minus $A$ in pp: means over splits, with the range covering 95\% of splits. AUC is for COVID-19 detection in $B$ / $A$.}
\label{tab:real}
\footnotesize
\begin{tabular}{@{}llcccccc@{}}
\toprule
Grouping & Scenario & Abstain $B$ / $A$ (\%) & Abstention gap & Silent-miss gap & Selective-risk gap & AUC $B$ / $A$ \\
\midrule
\multirow{5}{*}{\shortstack[l]{Sex\\(AUC 0.71)}} & No intervention & 23.3 / 18.5 & $+4.8$ [$-10.7$, $+21.7$] & $+6.9$ [$-4.4$, $+20.3$] & $+11.3$ [$-7.7$, $+31.1$] & 0.74 / 0.84 \\
 & Label bias & 22.7 / 19.0 & $+3.7$ [$-11.5$, $+21.9$] & $+12.9$ [$-3.3$, $+33.3$] & $+13.1$ [$-4.8$, $+33.8$] & 0.71 / 0.82 \\
 & View shift & 10.1 / 17.6 & $-7.6$ [$-20.2$, $+2.6$] & $-0.7$ [$-7.3$, $+5.6$] & $+8.7$ [$-1.6$, $+17.5$] & 0.59 / 0.83 \\
 & Scanner shift & \phantom{0}1.7 / 18.5 & $-16.9$ [$-31.1$, $-2.6$] & $+5.7$ [$-12.3$, $+92.0$] & $+27.6$ [$+2.9$, $+53.5$] & 0.45 / 0.84 \\
 & Label bias + scanner & \phantom{0}0.6 / 19.0 & $-18.4$ [$-30.7$, $-7.5$] & $+21.0$ [$-16.4$, $+96.0$] & $+27.8$ [$+3.6$, $+55.0$] & 0.48 / 0.82 \\
\midrule
\multirow{2}{*}{\shortstack[l]{Random\\(AUC 0.5)}} & No intervention & 19.8 / 20.5 & $-0.7$ [$-17.6$, $+15.9$] & $0.0$ [$-14.4$, $+11.8$] & $+0.5$ [$-21.7$, $+21.2$] & 0.80 / 0.80 \\
 & Label bias & 19.5 / 21.2 & $-1.7$ [$-16.4$, $+12.6$] & $0.0$ [$-17.0$, $+19.1$] & $+0.4$ [$-20.3$, $+23.8$] & 0.77 / 0.77 \\
\bottomrule
\end{tabular}
\end{table*}

The two label-bias rows reproduce the separability result on real images. Injected into female patients, whom the embeddings can partly identify, label bias raised the silent-miss gap by 6.0\pp{} relative to no intervention (an increase in 81\% of splits) and raised female patients' own silent-miss rate by 10.0\pp, while their abstention rate barely moved ($-0.6$\pp). Injected into a random group of the same size, the same bias raised that group's silent-miss rate by 6.8\pp{} but raised everyone else's by the same amount, so the gap stayed at zero: the harm was spread across all patients instead of being concentrated on one group.

Both real shifts lowered female patients' abstention, as in the simulation. The scanner shift cut it from 23.3\% to 1.7\% (lower in 99\% of splits) while discrimination collapsed below chance (AUC 0.45) and the error on accepted cases rose by 16.4\pp. In this case the harm appeared mainly as confident false positives, so the silent-miss gap did not grow, but the system stopped asking for help at the moment its predictions became least reliable. The view shift lowered abstention by 13.2\pp{} (lower in 97\% of splits) and discrimination by 0.14, without a measurable rise in silent misses or selective risk; supine images in this collection have a higher COVID-19 prevalence, which makes this comparison harder to interpret. Combined with label bias, the scanner shift produced the pattern of the simulation's worst case: abstention 18.4\pp{} lower for female patients and a silent-miss gap of +21.0\pp.
